\section{Von Neumann Model and ISA}

\begin{definition}
The \textbf{Von Neumann model} consists of memory (stores program and data), processing unit (ALU, registers), control unit (PC, IR), input, and output.
Key properties: \textbf{Stored Program} and \textbf{Sequential Processing}.
\end{definition}

\begin{remark}
Instruction Cycle: Fetch $\to$ Decode $\to$ Evaluate Address $\to$ Fetch Operands $\to$ Execute $\to$ Store Result.
\end{remark}

\begin{definition}
\textbf{Address Space}: Total unique memory locations.
\textbf{Addressability}: Bits stored per address.
\textbf{Endianness}: Order of bytes (Big Endian stores MSB at lowest address).
\end{definition}

\subsection{Instruction Set Architecture (ISA)}

\begin{definition}
The \textbf{ISA} is the interface between software and hardware. Defines memory organization, register set, and instructions (Opcode + Operands).
\end{definition}

\begin{lemma}
\textbf{CISC}: Complex instructions, small semantic gap.
\textbf{RISC}: Simple instructions, large semantic gap (e.g. MIPS, LC-3).
\end{lemma}

\begin{definition}
\textbf{Addressing Modes}: Immediate, Register, PC-Relative ($PC + offset$), Indirect, Base + Offset.
\end{definition}

\subsection{Assembly}

\begin{definition}
Assembly is human-readable machine code.
In RISC, operands must be in registers for ALUs:
\begin{itemize}
    \item \textbf{Load}: Reads data from memory into a register.
    \item \textbf{Store}: Writes data from a register into memory.
\end{itemize}
\end{definition}

\begin{remark}
Execution order altered by \textbf{Conditional Branches} (uses condition codes or register comparisons) and \textbf{Unconditional Jumps}.
\end{remark}
